\documentclass[aspectratio=169]{beamer}

\usetheme{Madrid}
\usecolortheme{default}

\title{Masterclass: Tax and Transfer Policy Evaluation}
\subtitle{Session 1: Structures In and Structures Out}
\author{Max Blesch}
\date{\today}

\begin{document}

\begin{frame}
  \titlepage
\end{frame}

\begin{frame}{Where to Find GETTSIM}
  \begin{itemize}
    \item Documentation (user-facing): \texttt{https://gettsim.readthedocs.io}
    \item Source code and issues: \texttt{https://github.com/ttsim-dev/gettsim}
    \item Personas package: \texttt{https://github.com/ttsim-dev/gettsim-personas}
  \end{itemize}
\end{frame}

\begin{frame}{GEPs: Why They Matter}
  \begin{itemize}
    \item GEPs are GETTSIM Enhancement Protocols: architecture and style rules.
    \item Key ones for users:
    \begin{itemize}
      \item GEP-07: user interface and \texttt{main()}.
      \item GEP-06: namespaces and qualified names.
      \item GEP-01/02: naming and data structure conventions.
      \item GEP-04: DAG-based computation model.
    \end{itemize}
    \item Use GEPs when API behavior seems non-obvious.
  \end{itemize}
\end{frame}

\begin{frame}{Session Workflow Around \texttt{main()}}
  \begin{itemize}
    \item Set \texttt{policy\_date\_str} first.
    \item Choose \texttt{main\_target} (what should come back).
    \item Provide \texttt{input\_data} (DataFrame+mapper or tree).
    \item Define \texttt{tt\_targets} (what to compute).
    \item Run and inspect outputs or missing-input diagnostics.
  \end{itemize}
\end{frame}

\begin{frame}{Diagnostic Workflow (Notebook)}
  \begin{itemize}
    \item Start with a small target set.
    \item Run once and read missing roots.
    \item Add coherent branches (not random leaves).
    \item Re-run and compare reduced missing lists.
    \item Use personas once the concept is clear.
  \end{itemize}
\end{frame}

\begin{frame}{What is \texttt{gettsim-personas}?}
  \begin{itemize}
    \item A collection of stylized household setups for GETTSIM.
    \item It provides coherent \texttt{input\_data\_tree} and often useful target bundles.
    \item In class: ideal for robust runs after learning custom DataFrame mapping.
  \end{itemize}
\end{frame}
\end{document}
